\documentclass{article}
\usepackage{graphicx} % Required for inserting images
\usepackage{amsmath}
\usepackage{amssymb}
\usepackage{enumerate}
\usepackage[margin=1in]{geometry}
\usepackage{tcolorbox}
\usepackage{bm}
\usepackage{graphicx}


\newcommand{\mydivider}{\vspace{1em}\hrule\vspace{1em}}
\newcommand{\mypart}[2]{\noindent{\textbf{[#1] point(s) --- #2:}}}
\newcommand{\programming}{\textcolor{blue}{This is a programming exercise. See \texttt{hw2.ipynb}}}



%%% BEGIN MACROS
% type your macros here
%%% END MACROS


\title{CS 577 --- Deep Learning --- Homework 2}
\author{}
\date{}

\begin{document}

\maketitle
\textbf{Read these instructions carefully}:
\begin{itemize}
  \item
In the \LaTeX~source code, type your answer in between
``\verb|%%% BEGIN ANSWER|'' and
``\verb|%%% END ANSWER|''.
        For advanced \LaTeX users,
        you can use your custom macros if you wish
        by placing them
        between
``\verb|%%% BEGIN MACROS|'' and
``\verb|%%% END MACROS|'' in the header.
        Do not modify anything else.




  \item Turn in both your \verb|.tex| file and the generated \verb|.pdf| file.
\end{itemize}



\section{Conditional Probability}
\mypart{3}{part a}
This exercise is designed to help you develop intuition for the concept of  conditional probability by ``creating one from scratch''.

In mathematics, a \emph{function} takes an input, let's call it \(x\), to exactly one output, let's call it \(y\).
For example, if
\({w} =
1/2
\)
and \(b = 1/2\), then the function
\begin{equation}
  \label{equation:function-with-description}
  y=
\underbrace{f({x})}_{\mbox{(name of the function)}} = \underbrace{w x + b}_{\mbox{(what the function actually is)}}
\end{equation}
assigns every input \(x\) to exactly one real number, namely \(wx + b\).
In this case, we say the  dependency of \(y\) on \(x\), i.e., the input-output relationship, is \emph{deterministic}.
Note in naming the function, we could have used \(g\) or \(h\) instead of \(f\).
There is nothing special about the letter ``\(f\)'' other than tradition (and also that the word ``function'' starts with that letter).

In the real world, it is often the case that
the dependency of \(y\) on \(x\) has a degree of uncertainty.
This notion is what \emph{conditional probability} aims to capture:
\begin{equation*}
  \underbrace{
    p(y\mid {x})
  }_{\mbox{(name of the conditional probability)}}
  =
  \underbrace{
    \qquad \qquad
  }_{\mbox{(what the conditional probability actually is)}}
  (\mbox{Intentionally blank})
\end{equation*}
Just like for functions, there is nothing special about the letter \(p\); we could've wrote \(q(y\mid \mathbf{x})\) or \(r(y \mid \mathbf{x})\).
Now, why did we leave the right-hand-side blank? We've learned how to ``create'' functions since grade school mathematics. But most of us never learned the notation for how to do the same for conditional probability.
Let's do that here.

Let \(n > 0\) be an integer.
Consider
\[y = wx+b + \epsilon\] where \(w\) and \(b\) are non-random numbers and \(x\) is the input, but \(\epsilon\) is random and defined via the following algorithm:
\begin{enumerate}
  \item Initialize \(\epsilon \gets 0\)

  \item For \(t = 1,\dots, n\), do:
        \begin{enumerate}[1.]
          \item Flip a fair coin

          \item If coin comes up head, then \(\epsilon \gets \epsilon + 1\)

          \item Else \(\epsilon \gets \epsilon - 1\)
        \end{enumerate}

  \item Return \(\epsilon/\sqrt{n}\).
\end{enumerate}
The above procedure defines a dependency of \(y\) on \(x\) that is random, and hence the input-output relationship is governed by a conditional probability.
Let's call the conditional probability
\[
  \mathcal{S}(y; wx+b, n)
\]
(the curly ``S'' is \verb|\mathcal{S}|).
When \(n=2\), we have
\[
  p(y\mid x) =
  \mathcal{S}(y; wx+b, 2)
  =
  \begin{cases}
    1/4
    &:
      \mbox{if }
      y =
   wx+b + \frac{2}{\sqrt{2}},
      \\
    1/2 &:
      \mbox{if }
          y = wx+b \\
    1/4
    &:
      \mbox{if }
      y=
   wx+b - \frac{2}{\sqrt{2}},
    \\
    0 &: \mbox{otherwise}.
  \end{cases}
\]
To see there, first note that there are \(4\) possible outcomes for the coin flip.
There is only one outcome in which both coins come up heads, then
\(
y=
   wx+b + \frac{2}{\sqrt{2}}
   \).
   Thus, the probability is \(1/4\). The same logic applies to
   the
\(
y=
   wx+b - \frac{2}{\sqrt{2}}
   \) case.
   Now, there are two outcomes that results in
   \(
    y = wx+b \): one coin comes up head while the other comes up tail. Hence the probability of this case is \(2/4 = 1/2\).
    It is impossible for \(y\) to equal to any other values.

    The exercise: calculate
\(  p(y\mid x) =
  \mathcal{S}(y; wx+b, 3)
\). \emph{Hint: introduce a notation for denoting outcomes of the coin flips.}


\begin{tcolorbox}
\textbf{Answer: }

%%% BEGIN ANSWER

\textbf{1:}
Each flip can result in either a head (H) or a tail (T). For \( n = 3 \), there are \( 2^3 = 8 \) possible outcomes:
\[
\text{HHH, HHT, HTH, HTT, THH, THT, TTH, TTT}.
\]

\textbf{2: }
The steps for determining \( \epsilon \) is:
- Start with \( \epsilon = 0 \),
- Add \( +1 \) for each head,
- Subtract \( -1 \) for each tail,
- Return \( \epsilon / \sqrt{3} \).

Let’s find \( \epsilon \) for each outcome:
\[
\begin{aligned}
    \text{HHH} &: \epsilon = 3 \quad \Rightarrow \quad \frac{3}{\sqrt{3}} = \sqrt{3}, \\
    \text{HHT} &: \epsilon = 1 \quad \Rightarrow \quad \frac{1}{\sqrt{3}} = \frac{1}{\sqrt{3}}, \\
    \text{HTH} &: \epsilon = 1 \quad \Rightarrow \quad \frac{1}{\sqrt{3}} = \frac{1}{\sqrt{3}}, \\
    \text{HTT} &: \epsilon = -1 \quad \Rightarrow \quad \frac{-1}{\sqrt{3}} = -\frac{1}{\sqrt{3}}, \\
    \text{THH} &: \epsilon = 1 \quad \Rightarrow \quad \frac{1}{\sqrt{3}} = \frac{1}{\sqrt{3}}, \\
    \text{THT} &: \epsilon = -1 \quad \Rightarrow \quad \frac{-1}{\sqrt{3}} = -\frac{1}{\sqrt{3}}, \\
    \text{TTH} &: \epsilon = -1 \quad \Rightarrow \quad \frac{-1}{\sqrt{3}} = -\frac{1}{\sqrt{3}}, \\
    \text{TTT} &: \epsilon = -3 \quad \Rightarrow \quad \frac{-3}{\sqrt{3}} = -\sqrt{3}.
\end{aligned}
\]

\textbf{}
Using the formula \( y = wx + b + \epsilon \), we get the possible values of \( y \) based on \( \epsilon \):
\[
\begin{aligned}
    y &= wx + b + \sqrt{3}, \\
    y &= wx + b + \frac{1}{\sqrt{3}}, \\
    y &= wx + b - \frac{1}{\sqrt{3}}, \\
    y &= wx + b - \sqrt{3}.
\end{aligned}
\]

\textbf{4:}
There are 8 possible outcomes from the coin flips, and the probabilities are:
\[
\begin{aligned}
    \mathbb{P}(y = wx + b + \sqrt{3}) &= \frac{1}{8} \quad \text{(HHH)}, \\
    \mathbb{P}(y = wx + b + \frac{1}{\sqrt{3}}) &= \frac{3}{8} \quad \text{(HHT, HTH, THH)}, \\
    \mathbb{P}(y = wx + b - \frac{1}{\sqrt{3}}) &= \frac{3}{8} \quad \text{(HTT, THT, TTH)}, \\
    \mathbb{P}(y = wx + b - \sqrt{3}) &= \frac{1}{8} \quad \text{(TTT)}.
\end{aligned}
\]

\textbf{Final Answer:}
Thus, the conditional probability \( p(y | x) = \mathcal{S}(y; wx + b, 3) \) is:
\[
p(y | x) = \mathcal{S}(y; wx + b, 3) =
\begin{cases}
    \frac{1}{8}, & \text{if } y = wx + b + \sqrt{3}, \\
    \frac{3}{8}, & \text{if } y = wx + b + \frac{1}{\sqrt{3}}, \\
    \frac{3}{8}, & \text{if } y = wx + b - \frac{1}{\sqrt{3}}, \\
    \frac{1}{8}, & \text{if } y = wx + b - \sqrt{3}, \\
    0, & \text{otherwise}.
\end{cases}
\]
%%% END ANSWER


\end{tcolorbox}


\mydivider

\mypart{2}{part b} \programming

\mydivider

\mypart{2}{part c} \programming

\section{The Gaussian distribution}



\mypart{1}{part a} Consider linear regression when the data live in the \(2\)-dimensional space. That is, the data vectors \(\mathbf{x}^{(i)} \in \mathbb{R}^{2}\) are \(2\)-dimensional vectors, and the labels \(y \in \mathbb{R}\) are numbers.



In lecture 2, we saw that for linear regression, the relationship between the sample (input) \(\mathbf{x}\) and the label (output) \(y\) is governed by
\begin{equation}
  \label{equation:conditional-distribution-for-linear-regression}
  p(y\mid \mathbf{x})
  =
  \mathcal{N}( y ;
  \mathbf{w}^{\top} \mathbf{x} +b , \, \sigma^{2}
  )
\end{equation}

where
  \(\mathcal{N}(\, \cdot \, ; \mu,\sigma^{2})
  \)
  is the probability density function for a Gaussian random variable with mean \(\mu\) and variance \(\sigma^{2}\).
The formula for
  \(\mathcal{N}(\, \cdot \, ; \mu,\sigma^{2})
  \)
  is given by
\[
  \mathcal{N}(z; \mu,\sigma^{2}) :=
\frac{1}{\sqrt{2\pi \sigma^{2}}}
\exp\left(
-\frac{1}{2} \frac{(z- \mu)^{2}}{\sigma^{2}}
\right).
\]

Let
\(\sigma = 1/\sqrt{2}\),
\(\mathbf{w} =
\begin{bmatrix}
0 \\ 1
\end{bmatrix}
\)
and \(b = 3\).
Compute \(p(y\mid \mathbf{x})\) when \(\mathbf{x} =
\begin{bmatrix}
  1 \\ -2
\end{bmatrix}
\).
The only unknown variable in your solution should be \(y\).
Type your answer in a single line of \LaTeX~code using fewer than 50 characters.
You don't need to show your work.


\begin{tcolorbox}
\textbf{Answer: }

%%% BEGIN ANSWER
$p(y \mid x) = \frac{1}{\sqrt{\pi}} \exp\left( -(y - 1)^2 \right).
$
%%% END ANSWER


\end{tcolorbox}


\mydivider

\mypart{1}{part b}
Plug in \(y = 1\) and \(y= -1/2\) into your answer above.
Compute the numerical value of \(\frac{p(1\mid \mathbf{x})} {p(-1/2\mid \mathbf{x})}\) and write the result below rounded to three decimal places (i.e., like 1.234)



\begin{tcolorbox}
\textbf{Answer: }

%%% BEGIN ANSWER
$\frac{p(1 \mid x)}{p\left(-\frac{1}{2} \mid x\right)} \approx 9.559

%%% END ANSWER

\end{tcolorbox}



\mydivider

\mypart{1}{part c} \programming

\mydivider

\mypart{2}{part d}
``\(p(y\mid \mathbf{x})= \mathcal{N}(y; \mathbf{w}^{\top} \mathbf{x} +b,\sigma^{2})\)''
is what is known as a probability \emph{density}:
\(p(y\mid \mathbf{x})\) is not the probability of observing exactly \(y\), but rather the probability
 of observing a value \emph{near} \(y\).
The number in the output cell of part c should be similar to the number in your answer in part b.
 Explain this why these two numbers are similar.



\begin{tcolorbox}
\textbf{Answer: }

%%% BEGIN ANSWER

The probability density function (PDF) $p(y \mid x) $ is like a guide that tells us how likely it is to see a value close to y. In part c, when we look at the ratio n1/n2, we're figuring out how often we see values of y within certain ranges. In part b, the ratio $ p(1 \mid x) / p(-1/2 \mid x) $ is about comparing how likely different specific values of y are.

Both of these ratios are similar because they both give us an idea of how likely it is to see certain values of y given x. The PDF gives us a theoretical view of this likelihood, while the empirical frequency shows us the actual data we’ve observed. In essence, both are ways to understand the same thing: how likely different values of y are when x is known.

%%% END ANSWER

\end{tcolorbox}

\section{A nonlinear regression problem}
\mypart{3}{part a}
Consider samples \(x \in \mathbb{R}\) and labels \(y \in \mathbb{R}\).
Let \(\bm{\theta} =
\begin{bmatrix}
  w_{1} \\ w_{2}
\end{bmatrix}
\)
and let \(f({x} ; \bm{\theta})
=
\sin(w_{1}x) + \cos(w_{2} x)
\).
Calculate the gradient \(\nabla_{\bm{\theta}} J(\bm{\theta})\) of the function
\[
  J(\bm{\theta})
  =
  J\left(
  \begin{bmatrix}
    w_{1}\\w_{2}
  \end{bmatrix}
\right  )^{2}
  =
  \frac{1}{N}
  \sum_{i=1}^{N} (y^{(i)} - f(\mathbf{x}^{(i)} ; \bm{\theta}))^{2}.
\]

\begin{tcolorbox}
\textbf{Answer: }

%%% BEGIN ANSWER


Lets compute the gradient $\nabla_\theta J(\theta)$ of the function 

\[
J(\theta) = \frac{1}{N} \sum_{i=1}^{N} \left( y^{(i)} - f(x^{(i)}; \theta) \right)^2,
\]
where 
\[
f(x; \theta) = \sin(w_1 x) + \cos(w_2 x),
\]
and $\theta = \begin{bmatrix} w_1 \\ w_2 \end{bmatrix}$.

\[\]

1. \textbf{}
Partial Derivatives of \(f(x; \theta)\) with respect to \(w_1\) and \(w_2\):
\[
\frac{\partial f(x; \theta)}{\partial w_1} = \frac{\partial}{\partial w_1} \left( \sin(w_1 x) + \cos(w_2 x) \right) = x \cos(w_1 x).
\]

\[
\frac{\partial f(x; \theta)}{\partial w_2} = \frac{\partial}{\partial w_2} \left( \sin(w_1 x) + \cos(w_2 x) \right) = -x \sin(w_2 x).
\]

2. \textbf{}
Applying the chain rule to get the gradient (\nabla_\theta J(\theta)):

\[
\nabla_\theta J(\theta) = \begin{bmatrix} \frac{\partial J(\theta)}{\partial w_1} \\ \frac{\partial J(\theta)}{\partial w_2} \end{bmatrix}.
\]

For the loss  \(J(\theta)\), we have:

\[
\frac{\partial J(\theta)}{\partial w_1} = -\frac{2}{N} \sum_{i=1}^{N} \left( y^{(i)} - f(x^{(i)}; \theta) \right) \frac{\partial f(x^{(i)}; \theta)}{\partial w_1},
\]
\[
\frac{\partial J(\theta)}{\partial w_2} = -\frac{2}{N} \sum_{i=1}^{N} \left( y^{(i)} - f(x^{(i)}; \theta) \right) \frac{\partial f(x^{(i)}; \theta)}{\partial w_2}.
\]

On Substituting the partial derivatives:

\[
\frac{\partial J(\theta)}{\partial w_1} = -\frac{2}{N} \sum_{i=1}^{N} \left( y^{(i)} - f(x^{(i)}; \theta) \right) x^{(i)} \cos(w_1 x^{(i)}),
\]
\[
\frac{\partial J(\theta)}{\partial w_2} = -\frac{2}{N} \sum_{i=1}^{N} \left( y^{(i)} - f(x^{(i)}; \theta) \right) (-x^{(i)} \sin(w_2 x^{(i)})).
\]

This implies to:

\[
\frac{\partial J(\theta)}{\partial w_1} = -\frac{2}{N} \sum_{i=1}^{N} \left( y^{(i)} - f(x^{(i)}; \theta) \right) x^{(i)} \cos(w_1 x^{(i)}),
\]
\[
\frac{\partial J(\theta)}{\partial w_2} = \frac{2}{N} \sum_{i=1}^{N} \left( y^{(i)} - f(x^{(i)}; \theta) \right) x^{(i)} \sin(w_2 x^{(i)}).
\]

3. \textbf{}
The gradient \(\nabla_\theta J(\theta)\) is:

\[
\nabla_\theta J(\theta) = \begin{bmatrix} -\frac{2}{N} \sum_{i=1}^{N} \left( y^{(i)} - f(x^{(i)}; \theta) \right) x^{(i)} \cos(w_1 x^{(i)}) \\ \frac{2}{N} \sum_{i=1}^{N} \left( y^{(i)} - f(x^{(i)}; \theta) \right) x^{(i)} \sin(w_2 x^{(i)}) \end{bmatrix}.
\]



%%% END ANSWER

\end{tcolorbox}
\mydivider

\mypart{1}{part b} \programming

\mydivider

\mypart{4}{part c}
Run gradient descent for \(100\) steps with constant step size \(=1/4\). Compare two different initialization: initialization ``A''
\[
w_{1} = -2.5,\quad
w_{2} = 3.56
\]
and initialization ``B''
\[
w_{1} = -2.5,\quad
w_{2} = 3.57.
\]
Answer the following questions:
\begin{enumerate}
  \item Which initialization led to better performance in terms of the training error?

  \item Is one of the converged final parameter (nearly) optimal? How do you know?
\end{enumerate}

\textcolor{red}{For this part, only your response below will be graded. We will not take into consideration of the contents in the python notebook.}
\begin{tcolorbox}
\textbf{Answer: }

%%% BEGIN ANSWER

1. \textbf{Which initialization led to better performance in terms of the training error?} \\
\[\]
Initialization B performed better, achieving a much lower final MSE of 0.0075 compared to Initialization A's final MSE of 1.0072. This indicates a more effective convergence.
\[\]
2. \textbf{Is one of the converged final parameters (nearly) optimal? How do you know?} \\
\[\]
Yes, the final parameters from Initialization B appear nearly optimal. The MSE is very close to zero, suggesting that the model parameters have converged to a point minimizing the training error effectively.

%%% END ANSWER

\end{tcolorbox}

\mydivider

\mypart{10 bonus}{part d} Explain in as much quantitative detail as possible why these two initialization led to different results. If you wish to use figures, you may use

\verb|\includegraphics[width=0.5\textwidth]{myfigure.png}|

Do not submit the figures separately.

\textcolor{red}{For this part, only your response below will be graded. We will not take into consideration of the contents in the python notebook.}

\begin{tcolorbox}
\textbf{Answer: }

%%% BEGIN ANSWER
\textbf

\textbf{Quantitative Explanation for Different Results:}

   1. Initialization A converged to a final MSE of 1.0072, whereas Initialization B achieved a much lower final MSE of 0.0075. This indicates that Initialization B found a more optimal set of parameters.

   2. Between iterations 10 and 20, Initialization B showed a sharp decrease in MSE from 1.1119 to 0.1855, while Initialization A's MSE only decreased marginally from 1.1419 to 1.0073 in the same range.

   3. The steeper descent in Initialization B suggests it started closer to a region with a more favorable gradient, leading to faster convergence. Initialization A plateaued earlier, showing less improvement over iterations.

   4. By iteration 25, Initialization B had already reached an MSE of 0.0075, while Initialization A was still at 1.0072, indicating that B found a better local minimum much faster.
\[\]
\includegraphics[width=\linewidth]{image1.png}


%%% END ANSWER
\end{tcolorbox}


\end{document}
